\begin{definition}[Face trace]
\label{def:face-trace}
For $\eta\in\Omega^{k-1}_{\mathrm{sm}}(I^k)$ and
$1\le j\le k$, $\varepsilon\in\{0,1\}$, set
\[
F_{j,\varepsilon}:=\{x\in I^k:\ x_j=\varepsilon\},
\qquad
\iota_{j,\varepsilon}:F_{j,\varepsilon}\hookrightarrow I^k .
\]
The full codimension-one face trace is the ordered tuple
\[
\Tr_{\partial I^k}^{\mathrm{face}}(\eta)
:=
\bigl(\iota_{j,\varepsilon}^\ast\eta\bigr)_{
1\le j\le k,\ \varepsilon\in\{0,1\}}.
\]
On $F_{j,\varepsilon}$ we use the coordinate orientation
\[
\dd x_1\wedge\cdots\wedge\widehat{\dd x_j}\wedge\cdots\wedge\dd x_k.
\]
The outward-oriented boundary sign is
\[
s_{j,1}:=(-1)^{j-1},
\qquad
s_{j,0}:=-(-1)^{j-1},
\]
so that
\[
\partial I^k=
\sum_{j=1}^k\sum_{\varepsilon\in\{0,1\}}
s_{j,\varepsilon}F_{j,\varepsilon}
\]
with the above coordinate face orientations.
\end{definition}

\begin{theorem}[Boundary determination on the unit cube]
\label{thm:boundary-rigidity}
Let $\eta\in\Omega^{k-1}_{\mathrm{sm}}(I^k)$ satisfy
$\dd\eta=\omega_0:=\dd x_1\wedge\cdots\wedge\dd x_k$. Then
\[
\norm{\eta}_{\coeff,\infty}\ge \frac1{2k}.
\]
If equality holds, then the full codimension-one face trace is forced: for every
$1\le j\le k$ and $\varepsilon\in\{0,1\}$,
\[
\iota_{j,\varepsilon}^\ast\eta
=
(-1)^{j-1}\frac{\varepsilon-\frac12}{k}\,
\dd x_1\wedge\cdots\wedge\widehat{\dd x_j}\wedge\cdots\wedge\dd x_k .
\]
Equivalently, the boundary trace of every sharp unit-cube minimizer is the
trace of the affine minimizer in Theorem~\ref{thm:box-minimizer-structure},
with no smaller coefficient-sup primitive possible.
\end{theorem}

\begin{proof}
Apply Theorem~\ref{thm:box-boundary-stability} with $R=I^k$. Then
$m_R=1/(2k)$, so its lower bound gives
$\norm{\eta}_{\coeff,\infty}\ge1/(2k)$. If equality holds, then $\eta$ is a
minimizer. The final assertion of Theorem~\ref{thm:box-boundary-stability}
therefore gives
\[
\iota_{j,\varepsilon}^\ast\eta
=s_{j,\varepsilon}\frac{1}{2k}\,
\dd x_1\wedge\cdots\wedge\widehat{\dd x_j}\wedge\cdots\wedge\dd x_k .
\]
The identity
$s_{j,\varepsilon}/(2k)=(-1)^{j-1}(\varepsilon-\frac12)/k$ is exactly the
displayed formula. The equivalence with the affine-minimizer trace is the
specialization of the minimizer-kernel description in
Theorem~\ref{thm:box-minimizer-structure} to $R=I^k$.
\end{proof}

\begin{corollary}[Quantitative trace stability on the unit cube]
\label{cor:cube-boundary-stability}
Let $\eta\in\Omega^{k-1}_{\mathrm{sm}}(I^k)$ satisfy
$\dd\eta=\omega_0:=\dd x_1\wedge\cdots\wedge\dd x_k$, and put
$M:=\norm{\eta}_{\coeff,\infty}$. Then $M\ge 1/(2k)$ and
\begin{equation}
\label{eq:cube-boundary-trace-stability}
\sum_{j=1}^k\sum_{\varepsilon\in\{0,1\}}
\left\|
\iota_{j,\varepsilon}^\ast\eta
-
(-1)^{j-1}\frac{\varepsilon-\frac12}{k}\,
\dd x_1\wedge\cdots\wedge\widehat{\dd x_j}\wedge\cdots\wedge\dd x_k
\right\|_{L^1(F_{j,\varepsilon};\coeff)}
\le
4k\left(M-\frac{1}{2k}\right).
\end{equation}
In particular, the exact boundary formula of
Theorem~\ref{thm:boundary-rigidity} follows when $M=1/(2k)$.
\end{corollary}

\begin{proof}
Apply Theorem~\ref{thm:box-boundary-stability} with $R=I^k$. Then
$m_R=1/(2k)$ and $P_1(I^k)=2k$. The affine minimizer has the displayed face
trace, so \eqref{eq:box-boundary-trace-stability} becomes
\eqref{eq:cube-boundary-trace-stability}. The lower bound $M\ge1/(2k)$ and the
equality case are also part of Theorem~\ref{thm:box-boundary-stability}.
\end{proof}

\begin{corollary}[Deficit form of cube trace stability]
\label{cor:cube-delta-trace-stability}
Let $\eta\in\Omega^{k-1}_{\mathrm{sm}}(I^k)$ satisfy
$\dd\eta=\omega_0:=\dd x_1\wedge\cdots\wedge\dd x_k$. If, for some
$\delta\ge0$,
\[
\norm{\eta}_{\coeff,\infty}\le \frac{1}{2k}+\delta,
\]
then
\begin{equation}
\label{eq:cube-delta-trace-stability}
\sum_{j=1}^k\sum_{\varepsilon\in\{0,1\}}
\left\|
\iota_{j,\varepsilon}^\ast\eta
-
(-1)^{j-1}\frac{\varepsilon-\frac12}{k}\,
\dd x_1\wedge\cdots\wedge\widehat{\dd x_j}\wedge\cdots\wedge\dd x_k
\right\|_{L^1(F_{j,\varepsilon};\coeff)}
\le
4k\delta .
\end{equation}
In particular, when $\delta=0$ this is exactly the boundary determination
formula in Theorem~\ref{thm:boundary-rigidity}.
\end{corollary}

\begin{proof}
Apply Corollary~\ref{cor:cube-boundary-stability}. With
$M:=\norm{\eta}_{\coeff,\infty}$, its right-hand side is
$4k(M-1/(2k))$. The hypotheses give $0\le M-1/(2k)\le\delta$, where the lower
bound is also part of Corollary~\ref{cor:cube-boundary-stability}. This proves
\eqref{eq:cube-delta-trace-stability}, and the final sentence follows by
setting $\delta=0$.
\end{proof}

\begin{corollary}[The centered radial primitive is extremal]
\label{cor:explicit-minimizer}
For $\omega_0=\dd x_1\wedge\cdots\wedge\dd x_k$ on $I^k$, the centered
radial primitive $K_k\omega_0$ equals the affine minimizer $\eta_R$ from
Theorem~\ref{thm:box-boundary-stability} with $R=I^k$. Hence
\[
\dd K_k\omega_0=\omega_0,
\qquad
\norm{K_k\omega_0}_{\coeff,\infty}=\frac{1}{2k},
\]
and $K_k\omega_0$ has the canonical face trace in
Theorem~\ref{thm:boundary-rigidity}.
\end{corollary}

\begin{proof}
For $R=I^k$ one has $m_R=1/(2k)$, so the formula for $\eta_R$ becomes
\[
\eta_R
=
\frac1k\sum_{j=1}^k(-1)^{j-1}\left(x_j-\frac12\right)
\dd x_1\wedge\cdots\wedge\widehat{\dd x_j}\wedge\cdots\wedge\dd x_k,
\]
which is exactly \eqref{eq:explicit-minimizer}. The norm, exterior derivative,
and trace statements follow from Theorem~\ref{thm:box-boundary-stability} and
Theorem~\ref{thm:boundary-rigidity}.
\end{proof}

\begin{theorem}[Sliced boundary readout for constant coordinate-monomial data]
\label{thm:sliced-active-readout}
Let
\[
Q:=\prod_{a=1}^n[0,L_a],
\qquad
I=\{i_1<\cdots<i_k\}\subset\{1,\dots,n\},
\qquad
\omega_I:=\dd x_{i_1}\wedge\cdots\wedge\dd x_{i_k},
\]
and set
\[
m_I:=\frac{1}{2\sum_{r=1}^kL_{i_r}^{-1}},
\qquad
Q_{I^c}:=\prod_{a\notin I}[0,L_a].
\]
Let $\eta\in\Omega^{k-1}_{\mathrm{sm}}(Q)$ satisfy $\dd\eta=\omega_I$, and
put $M:=\norm{\eta}_{\coeff,\infty}$. For
$1\le r\le k$ and $\lambda\in\{0,L_{i_r}\}$ define the active face
\[
G_{r,\lambda}:=\{x\in Q:\ x_{i_r}=\lambda\}
\]
and let $f_{r,\lambda}$ be the coefficient of
$\dd x_{I\setminus\{i_r\}}$ in $\iota_{r,\lambda}^\ast\eta$, that is,
\[
\coeff_{I\setminus\{i_r\}}\bigl(\iota_{r,\lambda}^\ast\eta\bigr)
=f_{r,\lambda}.
\]
With
\[
s_{r,L_{i_r}}:=(-1)^{r-1},
\qquad
s_{r,0}:=-(-1)^{r-1},
\]
one has $M\ge m_I$ and
\begin{equation}
\label{eq:sliced-active-stability}
\sum_{r=1}^k\sum_{\lambda\in\{0,L_{i_r}\}}
\int_{G_{r,\lambda}}
\abs{f_{r,\lambda}-s_{r,\lambda}m_I}\,\dd\Haus^{n-1}
\le
2\,\Vol(Q_{I^c})\,P_1(Q_I)(M-m_I),
\end{equation}
where $Q_I:=\prod_{r=1}^k[0,L_{i_r}]$. In particular, if $M=m_I$, then the
active coefficient on every active face is fixed by
\[
f_{r,\lambda}=s_{r,\lambda}m_I
\qquad
\text{for all }r,\lambda.
\]
No assertion is made about face components involving inactive differentials.
\end{theorem}

\begin{proof}
Fix the inactive variables $z=(x_a)_{a\notin I}\in Q_{I^c}$ and let
\[
\sigma_z:Q_I\to Q
\]
be the active slice map. Since $\dd\eta=\omega_I$, the pulled-back
$(k-1)$-form $\sigma_z^\ast\eta$ satisfies
\[
\dd(\sigma_z^\ast\eta)=\dd x_{i_1}\wedge\cdots\wedge\dd x_{i_k}
\]
on the active box $Q_I$. Its coefficient norm is at most $M$. Applying
Theorem~\ref{thm:box-boundary-stability} on this slice gives
\[
\sum_{r=1}^k\sum_{\lambda\in\{0,L_{i_r}\}}
\int_{G_{r,\lambda}\cap\sigma_z(Q_I)}
\abs{f_{r,\lambda}-s_{r,\lambda}m_I}\,\dd\Haus^{k-1}
\le
2P_1(Q_I)(M-m_I).
\]
The same slice theorem also gives $M\ge m_I$. Integrating the displayed
inequality over $z\in Q_{I^c}$ and using Fubini proves
\eqref{eq:sliced-active-stability}. Pullback to an active slice kills every
component containing an inactive differential, so the argument controls only
the coefficient of $\dd x_{I\setminus\{i_r\}}$ on the active faces.
\end{proof}

\begin{proposition}[Sharp invisibility of inactive face components]
\label{prop:inactive-face-invisibility}
Assume in addition that $n>k\ge2$. Let
\[
\eta_I:=2m_I\sum_{r=1}^k(-1)^{r-1}
\frac{x_{i_r}-L_{i_r}/2}{L_{i_r}}\,
\dd x_{i_1}\wedge\cdots\wedge\widehat{\dd x_{i_r}}\wedge\cdots
\wedge\dd x_{i_k}
\]
be the active affine primitive on $Q$. Let $J\subset\{1,\dots,n\}$ have
$|J|=k-1$ and contain at least one inactive index. For every constant
$c\in\RR$ with $|c|\le m_I$, set
\[
\eta_{J,c}:=\eta_I+c\,\dd x_J .
\]
Then $\dd\eta_{J,c}=\omega_I$,
$\norm{\eta_{J,c}}_{\coeff,\infty}=m_I$, and every active coefficient
controlled in Theorem~\ref{thm:sliced-active-readout} is identical for
$\eta_{J,c}$ and $\eta_I$. If $i_r\notin J$, however, the pullback to each
active face $G_{r,\lambda}$ contains the additional inactive component
$c\,\dd x_J$. Consequently exact norm saturation and the active readout data do
not determine inactive face components; the restricted conclusion of
Theorem~\ref{thm:sliced-active-readout} is sharp.
\end{proposition}

\begin{proof}
The form $\dd x_J$ is constant and closed, and $J$ contains an inactive index.
Therefore it is distinct from every active basis form
$\dd x_{I\setminus\{i_r\}}$, and
\[
\dd\eta_{J,c}=\dd\eta_I=\omega_I .
\]
The active affine primitive has coefficient norm $m_I$. Adding the new constant
coefficient $c$ in a distinct coordinate direction changes the coefficient
supremum to $\max\{m_I,|c|\}=m_I$.

On an active face $G_{r,\lambda}$, the coefficient measured in
Theorem~\ref{thm:sliced-active-readout} is the coefficient of
$\dd x_{I\setminus\{i_r\}}$. Since $J$ contains an inactive index, the added
term $c\,\dd x_J$ does not contribute to that coefficient, so all active
readout data are unchanged. If $i_r\notin J$, the pullback of $\dd x_J$ to
$G_{r,\lambda}$ is still $\dd x_J$, now regarded as a coordinate form on the
face, and hence the inactive face component has value $c$. Varying $c$ in the
interval $[-m_I,m_I]$ changes this component while preserving the equation,
the sharp norm, and all active readout coefficients.
\end{proof}

\begin{remark}
When $n>k$, Theorem~\ref{thm:sliced-active-readout} determines only the active
coefficient $\coeff_{I\setminus\{i_r\}}$ on the active faces
$G_{r,\varepsilon}$. Components involving inactive differentials are invisible
to the $k$-dimensional slice problem; Proposition~\ref{prop:inactive-face-invisibility}
shows that this invisibility persists even at exact saturation. The theorem
therefore identifies the maximal boundary information that is rigid under the
present inverse data.
\end{remark}

\begin{remark}
For $k\ge 2$, Theorem~\ref{thm:boundary-rigidity} yields rigidity of the
boundary trace, but not uniqueness in the interior. If
$\psi$ is a smooth $(k-2)$-form compactly supported in the interior of $I^k$
and $\eta_\varepsilon:=\eta_\ast+\varepsilon\,\dd\psi$, then
$\dd\eta_\varepsilon=\omega_0$ for every $\varepsilon$,
$\Tr_{\partial I^k}(\eta_\varepsilon)=\Tr_{\partial I^k}(\eta_\ast)$, and
$\norm{\eta_\varepsilon}_{\coeff,\infty}=1/(2k)$ for all sufficiently small
$|\varepsilon|$. Indeed, $\dd\psi$ vanishes on a neighborhood of
$\partial I^k$, where $\eta_\ast$ attains the value $1/(2k)$ in the
coefficient norm, whereas $\abs{\eta_\ast}_{\coeff,\infty}<1/(2k)$ on the
compact support of $\dd\psi$.
\end{remark}

